%%%%%%%%%%%%%%%%%%%%%%%%%%%%%%%%%%%%%%%
\chapter{Background and related work}
\label{chap:background}
%%%%%%%%%%%%%%%%%%%%%%%%%%%%%%%%%%%%%%%

L'état de l'art constitue un état des connaissances existantes, à un moment donné, sur un objet d'étude. Il s'agit donc de consulter la littérature scientifique relative à votre sujet et d'en faire une synthèse. Il ne s'agit pas de faire un \og catalogue \fg{} en juxtaposant et compilant les informations retenues. En rédigeant l'état de l'art, votre objectif est de justifier votre problématique, de démontrer sa pertinence et son intérêt au regard du contexte actuel.

Toujours prévoir un paragraphe avant une section. Une présentation générale, d'une ou deux lignes peuvent suffire.

\section{Titre de section}

Toujours prévoir un paragraphe avant une section. Une présentation
générale, d'une ou deux lignes peuvent suffire. Reprendre l'objet du
titre. Voir exemple suivant.

\subsection{Télétransportation quantique sur IP}

La télétransportation quantique sur le protocole IP\footnote{Ce
  protocole a été inventé par le Prof. Chacal au V\ieme siècle, lors de
  l'invasion des Nuls par les Huns.} permet de transférer
instantanément un corps à distance par la conversion d'unités
quantiques en trames IP. Le résultat s'avère cependant hasardeux
lorsque des trames se perdent.

\subsection{Titre de subsection}

Toujours prévoir un paragraphe avant une section. Une présentation
générale, d'une ou deux lignes peuvent suffire. Reprendre l'objet du
titre.


\subsubsection{Titre de subsubsection}

\section{Synthèse}

Il peut être utile pour certains chapitres complexes d'ajouter une
section \og Synthèse\fg{} qui résume le chapitre et ce qu'il faut en
retenir pour poursuivre sereinement la lecture des chapitres suivants.


\section{ Bibliographie ou Références}

Les références sur lesquelles se base un mémoire doivent être
scientifiques. Elles doivent être répertoriées de façon systématique,
complète et exacte. Toutes les références que vous reprenez dans votre
liste bibliographique doivent avoir été effectivement consultées. De
plus, chaque référence bibliographique doit correspondre au moins à un
‘renvoi' au sein de votre texte.

Il est recommandé d'utiliser l'outil \textbf{BibTeX} qui va faciliter
la gestion de votre bibliographie et attirera votre attention sur de
potentielles erreurs (omission de champs importants, etc).

Voici quelques références vers des entrées révélant les différentes subtilités lors de l'encodage: \cite{talon1972monologue},
\cite{toulemonde1968voisinage},
\cite{talon1975cuisine},
\cite{vanbroeck1980rhetorique},
\cite{crouton1983these},
\cite{talon1965memoire},
\cite{ministere1971rapport},
\cite{lefuneste1978solitude},
\cite{anon1969voisin},
\cite{talon2001blog},
\cite{talon1990oeuvrescompletes} et 
\cite{morose1985silence}.
